% ECONOMICS_PROBLEM: {"id": "E32-231", "title": "Self-fulfilling fluctuations", "jel_code": "E32", "source_id": "OP-0235"}
% !TeX program = xelatex
\documentclass[11pt,a4paper]{article}
\usepackage[margin=23mm]{geometry}
\usepackage{fontspec}
\setmainfont{Latin Modern Roman}
\usepackage{amsmath,amssymb,mathtools}
\usepackage{xcolor,enumitem,booktabs,longtable,array,xurl,hyperref}
\definecolor{muted}{HTML}{53636C}
\hypersetup{colorlinks=true,urlcolor=blue,linkcolor=blue}
\setlength{\parindent}{0pt}
\setlength{\parskip}{5pt}
\newcommand{\R}{\mathbb R}
\newcommand{\E}{\mathbb E}
\newcommand{\N}{\mathbb N}
\newcommand{\ind}{\mathbf 1}
\newcommand{\Var}{\operatorname{Var}}
\newcommand{\Cov}{\operatorname{Cov}}
\newcommand{\supp}{\operatorname{supp}}
\DeclareMathOperator*{\argmin}{arg\,min}
\DeclareMathOperator*{\argmax}{arg\,max}
\newcommand{\entrymeta}[3]{{\sffamily\small\color{muted}\textbf{\texttt{#1}}\quad|\quad #2\par #3\par}}
\newcommand{\Problemref}[1]{\texttt{#1}}
\newcommand{\JELref}[2]{\texttt{#2} (JEL \texttt{#1})}
\begin{document}
\section*{Self-fulfilling fluctuations}
\textbf{Problem ID:} \texttt{E32-231}\quad \textbf{JEL:} \texttt{E32}\par
% BEGIN ECONOMICS STATEMENT
\label{problem:OP-0235}
{\sffamily\small\color{muted}Original subject group: Business cycles and macro dynamics\par}
\label{definitions:problem:30007516}
\textbf{Audit classification:} Explicit published open question. A quantitative bubbles-and-crashes challenge is explicit; the cross-world recession attribution is editorial. Verification concerns the cited version, not first priority or proof that this exact editorial specification remains unsolved on October 8, 2026.
\subsection*{Definitions and precise research target}
\paragraph{Editorial mathematical specification.} Fix a nonlinear incomplete-markets model with observed fundamental state \(f_t\), endogenous distribution \(\mu_t\), and a declared equilibrium-selection variable \(b_t\) representing beliefs or coordination. A belief innovation \(\eta_t\) is required to be orthogonal to innovations in the specified fundamental-information set, not merely to lagged output. Let equilibrium log output be \(y_t=g(f_t,\mu_t,b_t)\). Define a recession as \(y_{t+H}-y_t<-q\), for preselected horizon \(H\) and threshold \(q>0\).
For each observed recession episode, construct the counterfactual path \(y^{0}_t\) obtained by suppressing \(\eta_t\) while retaining the same fundamental innovations, initial distribution, and policy rule. The target is
\[
\Theta=\Pr\{y_{t+H}-y_t<-q,\;
y^{0}_{t+H}-y^{0}_t\ge -q\},
\]
the probability that a recession would be prevented under this intervention on beliefs. Determine an identified set for \(\Theta\) from forecast surveys, balance sheets, asset prices, and instruments with explicit exclusion restrictions. Because equilibrium selection is latent and beliefs can react to fundamentals, provide observationally equivalent explanations and falsifiable restrictions. Establish robustness to plausible selection rules and financial-friction specifications. Existence of sunspot equilibria in a particular model does not establish their quantitative importance in actual recessions.

A cross-world event needs a coupling. Specify measurable structural equations \(f_{t+1}=F(f_t,\xi_{t+1})\), \(b_{t+1}=B(b_t,f_t,\mu_t,\eta_{t+1})\), and \(\mu_{t+1}=U(f_t,\mu_t,b_t)\), together with an initial law and joint law of \((\xi,\eta)\). Both worlds use the same initial draw and fundamental innovations; the intervention replaces the specified future \(\eta\) sequence by zero and recomputes \(b,\mu,y\), rather than holding the endogenous distribution fixed. A sunspot is extrinsic randomness absent from preferences, technology, endowments, and information about fundamentals; a nonfundamental belief shock need not be a rational-expectations sunspot. Let \(D\) and \(D^0\) be the two recession indicators. Then \(\Theta=P(D=1,D^0=0)\), while the fraction of observed recessions prevented is \(\Theta/P(D=1)\) when \(P(D=1)>0\). Marginal recession probabilities alone do not identify their joint law.
\subsection*{Variants and assumption changes}
The following are \emph{editorial variants}, not additional published open conjectures.
\begin{enumerate}
\item \emph{No cross-world coupling.} Identify only \(p=P(D=1)\) and \(p_0=P(D^0=1)\). Sharp unrestricted coupling bounds are \(\max(0,p-p_0)\le\Theta\le\min(p,1-p_0)\); compare these with bounds under a specified monotonicity assumption.
\item \emph{Rational sunspot equilibrium.} Require agents' conditional forecasts to coincide with the selected equilibrium law, and let an independent extrinsic signal select equilibria. Determine whether quantitatively important crashes remain feasible under household portfolio moments; this excludes arbitrary subjective forecast errors.
\end{enumerate}
\subsection*{References and source audit}
\begin{itemize}
\item Benjamin Moll. \emph{Research Agenda: Benjamin Moll}. 2020. Locator: Section4.3, Heterogeneity, Bubbles and Crashes. Primary text checked. \url{https://economicdynamics.org/research-agenda-moll2020/}.
\end{itemize}
A quantitative bubbles-and-crashes challenge is explicit; the cross-world recession attribution is editorial.
\begin{enumerate}\raggedright
\item\hypertarget{definitions:source:30007516:1}{} Benjamin Moll. 2020. \emph{Research Agenda: Benjamin Moll}. Section4.3, Heterogeneity, Bubbles and Crashes.
\url{https://economicdynamics.org/research-agenda-moll2020/}.
\end{enumerate}

\subsection*{Common conventions: Source mathematical conventions}

These conventions apply to each entry unless explicitly replaced.

\textbf{Models and identification.} A primitive model \(m\) specifies all probability laws, feasible choices, information, equilibrium and measurement rules used by its target. The admissible class \(\mathcal M\) is fixed independently of outcomes. If \(P_O(m)\) is its observable probability law and \(\tau(m)\) the target, its identified set at observable law \(P\) is
\[
 \mathcal I_\tau(P)=\{\tau(m):m\in\mathcal M,\ P_O(m)=P\}.
\]
Point identification means that this set is a singleton. An empty set signals incompatibility of data and assumptions. Coordinate bounds need not describe the joint set. Bounds are sharp only if every retained target is attainable or arbitrarily approachable by compatible models. Finite bounds require explicit restrictions on unobserved quantities; unrestricted confounding, hidden wealth, extrapolation or arbitrary equilibrium selection can yield uninformative sets.

\textbf{Causal interventions.} A potential outcome \(Y_i(p)\) is the outcome under a complete policy regime \(p\), including timing, financing, information and market spillovers. The target population and units are fixed before treatment. With interference, use the complete assignment vector or a justified exposure mapping; individual no-interference notation alone cannot identify an economy-wide intervention. A difference of mean potential outcomes does not require the cross-world joint distribution. Conditioning on post-treatment migration, survival, enrollment or employment changes the estimand and needs separate justification.

\textbf{Statistical statements.} An estimator, confidence set, forecast or test specifies a sampling law, model class and loss/coverage criterion. Uniform \((1-\alpha)\)-coverage means \(\inf_{m\in\mathcal M}\Pr_m\{\tau(m)\in C_n\}\geq1-\alpha\), or an explicitly asymptotic version. The whole parameter space trivially has coverage, so informativeness requires a stated width, risk or power objective. A forecasting improvement is relative to a named benchmark, information set and evaluation population.

\textbf{Equilibrium and optimization.} An equilibrium comprises feasible plans, optimality, beliefs where needed, budget constraints and all market/resource clearing. The primitives or a stated benchmark define each correspondence \(\mathcal E(p;m)\). Nonemptiness is assumed or proved, never inferred just from compact policy sets. A selected-equilibrium target uses a declared selection rule; a robust target quantifies over all permitted equilibria. Use suprema or infima unless attainment is established. An \(\varepsilon\)-optimal policy is feasible and within \(\varepsilon\) of the finite optimum value. Report an empty feasible set separately.

\textbf{Welfare and accounting.} Normative utility, interpersonal normalization, social weights, numeraire, discounting and the use of public receipts are primitives. Behavioral utility need not coincide with the welfare criterion. Household welfare includes ownership income and rents once; adding profits again to compensating variation generally double-counts them. A monetary equivalent is defined by an explicit transfer and price convention, with monotonicity and range conditions. Average effects, mechanism contributions, transfers and welfare losses are different targets. Infinite sums require convergence and dynamic feasible plans require terminal or no-Ponzi conditions.

\textbf{Source versions.} A working-paper date, revision date and journal publication year are distinguished. A DOI for a journal version is not automatically the DOI of an earlier manuscript. Locators refer to the stated version and printed pages where specified. A metadata-only check does not verify a claim about a full-text conclusion. Every entry's source subsection narrows the provenance claim accordingly.
% END ECONOMICS STATEMENT
\end{document}
